\documentclass[conference]{IEEEtran}
\IEEEoverridecommandlockouts
\usepackage{cite}
\usepackage{amsmath,amssymb,amsfonts}
\usepackage{algorithmic}
\usepackage{graphicx}
\usepackage{grffile}
\usepackage{textcomp}
\usepackage{xcolor}
\usepackage{booktabs}
\usepackage{listings}
\usepackage{url}
\usepackage{tabularx}

% Make Abstract, Index Terms, Section, and Subsection headings Bold
\renewcommand{\abstractname}{\textbf{ABSTRACT}}
\renewcommand{\IEEEkeywordsname}{\textbf{Index Terms}}

\makeatletter
\renewcommand\section{\@startsection{section}{1}{\z@}%
  {1.5ex plus 1.5ex minus 0.5ex}%
  {0.7ex plus 0.5ex minus 0.2ex}%
  {\normalfont\normalsize\bfseries\centering}}
\renewcommand\subsection{\@startsection{subsection}{2}{\z@}%
  {1.2ex plus 1.0ex minus 0.2ex}%
  {0.5ex plus 0.2ex}%
  {\normalfont\normalsize\bfseries}}
\makeatother

\def\BibTeX{{\rm B\kern-.05em{\sc i\kern-.025em b}\kern-.08em
    T\kern-.1667em\lower.7ex\hbox{E}\kern-.125emX}}

\definecolor{codebg}{RGB}{248, 248, 248}
\definecolor{codeframe}{RGB}{210, 210, 210}
\lstset{
    backgroundcolor=\color{codebg},
    basicstyle=\ttfamily\tiny,
    breaklines=true,
    breakatwhitespace=false,
    frame=single,
    rulecolor=\color{codeframe},
    showstringspaces=false,
    keepspaces=true,
    columns=flexible,
    tabsize=2,
    extendedchars=true,
    literate={¦}{|}1 {+--}{|--}2 {+--}{`--}2 {-}{-}1
}

\begin{document}

\title{Microprocessor-Controlled Embedded NAS \\& Static Web Hosting Node\\
{\footnotesize \textsuperscript{*}Course Project Report: CSE0611324 -- Microprocessors, Microcontrollers and Assembly Language Lab}\\
{\footnotesize \textbf{Course Instructors / Evaluators:} Tahmina Yeasmin Lima (Lecturer) \\& Abrar Mahbub Tanim (Lecturer)}\\
{\footnotesize \textit{Dept. of Computer Science \\& Engineering, University of Information Technology and Sciences (UITS)}}}

\author{
\IEEEauthorblockN{Jannatul Ferdus}
\IEEEauthorblockA{\textit{Dept. of Computer Science \& Engineering} \\
\textit{University of Information Technology and Sciences (UITS)}\\
Dhaka, Bangladesh \\
0432410005101058}
\and[\linebreak\hfill]
\IEEEauthorblockN{Gaus Saraf Murady}
\IEEEauthorblockA{\textit{Dept. of Computer Science \& Engineering} \\
\textit{University of Information Technology and Sciences (UITS)}\\
Dhaka, Bangladesh \\
0432410005101088}
\and[\hfill\linebreak]
\IEEEauthorblockN{Sadia Islam}
\IEEEauthorblockA{\textit{Dept. of Computer Science \& Engineering} \\
\textit{University of Information Technology and Sciences (UITS)}\\
Dhaka, Bangladesh \\
0432410005101096}
}

\maketitle

\begin{abstract}
Traditional Network Attached Storage (NAS) units and web servers depend heavily on power-hungry x86 hardware or Linux single-board computers (SBCs) drawing between 15 and 60 Watts. In this work, we engineer and evaluate an ultra-low-power ($<1$ Watt), air-gapped embedded micro-NAS and static website hosting node powered by the Espressif ESP32 dual-core microcontroller. By coordinating FreeRTOS asymmetric multiprocessing, a dedicated hardware VSPI bus interface clocked at 4~MHz, and a fixed 2~KB streaming multipart buffer architecture, the node delivers browser-based file uploads, high-throughput downloads, storage telemetry, dynamic cluster reclamation, and static HTML5 web delivery from raw FAT32 block storage without cloud dependencies or client-side software. The design is rigorously validated within PlatformIO and a cycle-accurate local Wokwi simulation environment bridged via an open-socket Layer-2/3 gateway, proving multi-device interoperability across wired Ethernet and wireless hotspot subnets for under 750 Bangladeshi Taka (\$7 USD).
\end{abstract}

\begin{IEEEkeywords}
ESP32, Microprocessor Systems, Embedded NAS, FreeRTOS, VSPI Bus, FAT32 Storage, Static Web Server, Air-Gapped IoT.
\end{IEEEkeywords}

\section{\textbf{Introduction}}
Lightweight local data sharing and edge document distribution in field environments, disaster recovery zones, and academic laboratories rarely require the computational complexity, high thermal footprint, or power demand of general-purpose operating systems. Commercial NAS appliances and single-board computers such as the Raspberry Pi draw between 3 and 7 Watts under load and cost upwards of \$45 to \$75. Furthermore, managing full Linux installations introduces security vulnerabilities, unexpected kernel panics, and recurring software patch maintenance cycles.

In contrast, the ESP32 microcontroller costs approximately 400 BDT (\$4 USD) and consumes less than 0.5 Watts during standard operations. Armed with an Xtensa 32-bit dual-core microprocessor operating at 240~MHz and an integrated 802.11 b/g/n Wi-Fi radio, the ESP32 possesses sufficient hardware capabilities to handle direct block storage I/O and concurrent HTTP socket transactions.

This paper presents the comprehensive engineering, implementation, and empirical verification of ``Project.MM''---an autonomous, low-overhead micro-NAS and static web node. The system was developed using PlatformIO and cross-validated in a virtual simulation environment utilizing Wokwi and a custom TCP forwarder (\texttt{wokwigw}). By decoupling network communication from storage serialization across independent CPU cores, our firmware maintains a deterministic memory ceiling of approximately 2~KB RAM during high-volume multipart uploads and binary file downloads.

\section{\textbf{Project Overview \& Objectives}}

\subsection{\textbf{Problem Statement \& Architectural Challenges}}
Implementing a multi-client file server on a deeply constrained embedded platform faces three critical bottlenecks:
\begin{enumerate}
    \item \textbf{Severe Volatile Memory Limits:} The ESP32 features 520~KB of internal SRAM, with only roughly 300~KB available for dynamic heap allocation. Attempting to buffer entire incoming files in memory during HTTP multipart transfers triggers immediate out-of-memory kernel panics and watchdog resets.
    \item \textbf{Bus Contention \& Socket Starvation:} Standard Arduino library calls like \texttt{Stream::readBytes()} poll data character-by-character. When transferring multi-megabyte files across a serial peripheral interface, slow block writes block the CPU, preventing the network stack from processing TCP ACK packets and leading to socket deadlocks.
    \item \textbf{Virtual Simulation Bus Synchronization:} Cycle-accurate virtual simulators like Wokwi enforce strict SPI clocking and part naming schemas. Multi-sector write commands (\texttt{CMD25}) often fail under virtual bus latency unless decomposed into discrete single-sector transactions (\texttt{CMD24}).
\end{enumerate}

\subsection{\textbf{Core Objectives}}
To resolve these limitations, the system was designed around five core engineering objectives:
\begin{itemize}
    \item Build an autonomous, localized file server over the hardware VSPI bus supporting file uploads, downloads, and directory management.
    \item Develop an intuitive, mobile-responsive dark-mode user interface displaying real-time FAT32 capacity metrics (\texttt{SD.totalBytes()} vs. \texttt{SD.usedBytes()}) and recursive file indexing.
    \item Implement dual-purpose firmware capability: serving raw binary file downloads alongside hosted static HTML5/CSS web assets directly from the \texttt{/www/} directory.
    \item Enforce system integrity by applying read-only write protection to critical system folders (\texttt{/default}, \texttt{/www}) and dynamically seeding initial assets from Flash ROM via \texttt{PROGMEM}.
    \item Validate seamless multi-device connectivity across wired Ethernet, campus switches, and wireless mobile hotspots without requiring client software installation.
\end{itemize}

\section{\textbf{System Architecture \& Working Principle}}

\subsection{\textbf{Microprocessor \& Storage Interface}}
The system utilizes the classic ESP32-WROOM-32 module, containing an Xtensa Dual-Core 32-bit LX6 microprocessor with 520~KB internal SRAM and 4~MB external SPI Flash. Secondary block storage is provided by a standard MicroSD card breakout module formatted with FAT32.

To guarantee high signal integrity and prevent conflict with boot-critical strapping pins, the storage module is interfaced through the ESP32 hardware VSPI peripheral. Table~\ref{tab:pinout} specifies the exact pinout mapping.

\begin{table}[htbp]
\caption{Hardware VSPI to MicroSD Module Pinout}
\begin{center}
\begin{tabularx}{\columnwidth}{l l l X}
\toprule
\textbf{ESP32 Pin} & \textbf{SD Pin} & \textbf{Function} & \textbf{Description} \\
\midrule
3V3 & VCC & Power & Regulated 3.3V DC Supply \\
GND & GND & Ground & Common Ground Reference \\
GPIO 5 & CS & Chip Select & Active-Low Slave Select Line \\
GPIO 18 & SCK & Clock & SPI Clock (@ 4 MHz) \\
GPIO 23 & MOSI & Data In & Master Out, Slave In \\
GPIO 19 & MISO & Data Out & Master In, Slave Out \\
\bottomrule
\end{tabularx}
\label{tab:pinout}
\end{center}
\end{table}

\subsection{\textbf{Working Method \& Execution Flow}}
The firmware execution and runtime communications progress through a two-phase initialization and packet handling sequence, as illustrated in the system networking bridge topology in Fig.~\ref{fig:network_bridge}:
\begin{enumerate}
    \item \textbf{Phase 1 (Storage Mount \& Flash Asset Seeding):}
    The CPU initializes the SPI hardware lines via \texttt{SPI.begin(18, 19, 23, 5)} and mounts the FAT32 filesystem at 4~MHz (\texttt{SD.begin(5, SPI, 4000000)}). It creates standard folder trees (\texttt{/default}, \texttt{/www}, \texttt{/uploads}) and checks if default assets (\texttt{fileserver.html}, \texttt{index.html}, \texttt{lab\_report.txt}) exist on storage. If missing or 0~bytes, it dynamically seeds them from Flash ROM.
    \item \textbf{Phase 2 (Network Association \& HTTP Routing):}
    The node associates with the local wireless network using static channel binding (Channel 6) to eliminate multi-channel scan latency. Once connected, it binds REST routes (\texttt{/}, \texttt{/site}, \texttt{/api/files}, \texttt{/get}, \texttt{/upload}, \texttt{/delete}) and initializes the web daemon on port 80. Concurrently, a FreeRTOS diagnostic task verifies endpoint health via local loopback requests.
\end{enumerate}

\begin{figure}[htbp]
\centering
\includegraphics[width=0.58\columnwidth]{"Diagram/4.2 Networking Bridge Architecture (Wokwi Gateway).png"}
\caption{Layer-2/3 TCP forwarder and WebSocket bridging architecture connecting host browsers to the simulated ESP32 node via \texttt{wokwigw}.}
\label{fig:network_bridge}
\end{figure}

\subsection{\textbf{Strapping Pin Safety vs. SDMMC Trade-off}}
While the ESP32 supports a 4-bit parallel SDMMC interface capable of higher throughput (20--40~MB/s), SDMMC requires GPIO~12 (\texttt{MTDI}). GPIO~12 is an internal boot strapping pin; if pulled HIGH by an SD card pull-up resistor during power-up, the internal voltage regulator sets flash voltage to 1.8V instead of 3.3V, causing a permanent boot loop. In contrast, the dedicated VSPI bus on GPIO~5, 18, 19, and 23 completely avoids strapping pins, ensuring 100\% boot reliability across all hardware batches and simulation engines.

\subsection{\textbf{FreeRTOS Dual-Core Multiprocessing}}
The firmware exploits the Symmetric Multiprocessing (SMP) capabilities of FreeRTOS:
\begin{itemize}
    \item \textbf{Core 0 (Networking Engine):} Dedicated strictly to the LwIP (Lightweight IP) TCP/IP protocol stack, 802.11 radio drivers, and incoming packet interrupt handling.
    \item \textbf{Core 1 (Application \& Storage Engine):} Manages the primary web server loop (\texttt{server.handleClient()}), FAT32 cluster traversal, SPI sector read/write operations, and a pinned background diagnostic worker (\texttt{runClientTests}).
\end{itemize}
Because disk latency on Core 1 never preempts networking tasks on Core 0, client TCP connections maintain steady handshake keep-alives without dropouts.

\section{\textbf{Implementation Procedures \& Firmware Engineering}}

\subsection{\textbf{Zero-Overhead Chunked Streaming Architecture}}
To eliminate memory exhaustion, incoming multipart file uploads are parsed via \texttt{server.upload()} chunk-by-chunk. Each received network packet is flushed immediately into 512-byte sector-aligned SD buffers before accepting subsequent TCP windows. The runtime heap allocation remains locked at approximately 2~KB regardless of whether transferring a 10~KB assembly source file or a 20~MB binary archive.

Similarly, file downloads bypass standard buffered reads. Our custom \texttt{streamSDFile()} routine sets an internal buffer of 1,024 bytes and streams chunks directly to the active socket via non-blocking \texttt{client.write()} calls with explicit \texttt{vTaskDelay(1)} yields, completely preventing TCP buffer deadlocks.

\subsection{\textbf{Flash ROM PROGMEM Storage}}
To preserve internal SRAM, all production web templates (\texttt{fileserver.html} and \texttt{index.html}) are converted into raw byte arrays stored in Flash memory using the \texttt{PROGMEM} attribute (\texttt{default\_assets.h}). When client browsers request fallback content, templates are served directly from Flash ROM via \texttt{server.send\_P()}, saving over 300~KB of SRAM heap.

\section{\textbf{Low-Level Debugging Challenges \& Solutions}}

During development and simulation validation, several complex engineering challenges were identified and systematically resolved:

\subsection{\textbf{The Wokwi D-Pin Silkscreen Bug}}
In initial circuit descriptions (\texttt{diagram.json}), wiring was declared as \texttt{esp:5}. In physical hardware, GPIO~5 is pin 5; however, Wokwi's component schema strictly mandates the silkscreen identifier \texttt{esp:D5}. Declaring \texttt{esp:5} silently caused the CS line to remain floating, resulting in recurring \texttt{SD.begin()} initialization timeouts. Correcting the descriptor to \texttt{esp:D5} restored SPI bus communication.

\subsection{\textbf{Clock Schema Identifier Mismatch (\texttt{SCK} vs. \texttt{CLK})}}
Physical breakout boards frequently label the clock line as \texttt{CLK}. However, the virtual SD card schema in Wokwi mandates \texttt{sd1:SCK}. Connecting \texttt{esp:D18} to \texttt{sd1:CLK} left the serial clock unmapped, causing the virtual SPI state machine to freeze. Aligning the pin descriptor to \texttt{sd1:SCK} resolved all synchronization issues.

\subsection{\textbf{CMD24 Sector Write Decomposition}}
Virtual SD controllers in simulation frequently fail during multi-sector streaming writes (\texttt{CMD25}) when virtual latency deviates from physical silicon timing, resulting in 0-byte corrupt files. To resolve this, our upload routines decompose streaming writes into sequential, single-sector 512-byte blocks (\texttt{CMD24}) followed by explicit sector flushes. This guaranteed 100\% byte-accurate data commits across all simulated file operations.

\subsection{\textbf{High-Throughput Download Bottleneck Optimization}}
Default Arduino \texttt{WebServer::streamFile()} uses character-by-character polling via \texttt{Stream::readBytes()}. For files larger than 1~KB, this caused severe CPU thrashing and socket timeouts. Implementing \texttt{streamSDFile()} with a dedicated 1~KB transfer buffer and explicit socket flushing increased download throughput by over 1,200\% while completely eliminating socket lockups.

\section{\textbf{Network Architecture \& Laboratory Deployment Scenarios}}

\subsection{\textbf{Dual-Network Gateway Bridging (\texttt{wokwigw})}}
Virtual Wi-Fi inside the Wokwi engine operates within an isolated subnet (\texttt{10.13.37.2}). To make the simulated microcontroller accessible across physical devices, the companion bridge executable \texttt{wokwigw.exe} acts as a Layer-2/3 TCP forwarder:
\begin{itemize}
    \item Connects to Wokwi via WebSocket on \texttt{ws://localhost:9011}.
    \item Binds to \texttt{0.0.0.0:9080} (\texttt{INADDR\_ANY}) on the host machine.
    \item Routes inbound traffic on host port 9080 directly to virtual ESP32 port 80.
\end{itemize}
An automated batch script (\texttt{setup\_firewall.bat}) configures Windows Defender Firewall with auto-elevation to permit inbound TCP port 9080 traffic across Private, Domain, and Public network profiles.

\subsection{\textbf{The 4 Validated Simulation Scenarios}}
As detailed in our operational guides, the system was verified across four real-world network topologies summarized in Table~\ref{tab:scenarios}.

\begin{table}[htbp]
\caption{The Four Validated Network Simulation Scenarios}
\begin{center}
\begin{tabularx}{\columnwidth}{lXX}
\toprule
\textbf{Scenario} & \textbf{Network Topology} & \textbf{Key Operational Benefit} \\
\midrule
\textbf{1. Home Hybrid} & Host on wired Ethernet; client phones on Wi-Fi router. & Eliminates radio interference; highly stable high-volume testing. \\
\addlinespace
\textbf{2. Lab Switch} & Both Host PC and Client PC wired to campus lab switch. & \textbf{LocalSend Web Model:} Instant peer-to-peer 2-PC lab file sharing with zero install. \\
\addlinespace
\textbf{3. Mobile Hotspot} & Host laptop creates Windows Mobile Hotspot. & \textbf{Deterministic IP (192.168.137.1);} completely defeats campus Wi-Fi AP client isolation. \\
\addlinespace
\textbf{4. All-Wireless} & Host PC and client devices on common Wi-Fi SSID. & Standard home/dorm wireless access convenience. \\
\bottomrule
\end{tabularx}
\label{tab:scenarios}
\end{center}
\end{table}

\subsection{\textbf{LocalSend Web Model: Zero-Install 2-PC Lab Sharing}}
Scenario~2 replicates the popular LocalSend Web workflow. In computer laboratories, students frequently need to exchange assembly code files, test inputs, and lab reports between adjacent workstations without logging into personal email or cloud storage. With our system running on PC~1, student PC~2 simply opens Chrome or Edge to PC~1's IP address (\texttt{http://10.100.x.x:9080/}) and immediately uploads or downloads files directly from the ESP32 node.

\section{\textbf{Project Scope, Limitations \& Silicon Realities}}

\subsection{\textbf{Physical Decoupling Capacitors \& Brownout Prevention}}
In simulation, power rails are mathematically ideal. On physical hardware, Wi-Fi radio transmission bursts draw transient current spikes in excess of 300~mA. Concurrently, MicroSD NAND flash sector writes draw up to 100~mA. If powered over high-resistance breadboard jumper wires or weak USB ports, the voltage rail sags below the ESP32's 2.8V brownout threshold, causing spontaneous CPU resets. 

To transition from virtual simulation to physical breadboards, the circuit design mandates placing a 100~$\mu$F electrolytic capacitor alongside a 0.1~$\mu$F ceramic capacitor directly across the MicroSD module's VCC and GND pins.

\subsection{\textbf{SPI Bus Parasitic Capacitance \& Clock Capping}}
While simulation lines exhibit zero impedance, breadboard jumper wires longer than 15~cm introduce parasitic capacitance. Running the SPI clock at 20--40~MHz causes severe signal distortion and data corruption. Capping the bus clock frequency at 4~MHz to 10~MHz ensures pristine square wave integrity and reliable sector transfers.

\subsection{\textbf{HTTP vs. HTTPS Trade-Offs}}
The server operates over unencrypted HTTP rather than HTTPS. This is an intentional engineering trade-off: TLS 1.3 handshakes consume approximately 45~KB of volatile heap per connection and require trusted public Certificate Authority (CA) validation, which cannot be acquired for dynamic private IP addresses on local subnets. Operating on plaintext HTTP leaves over 280~KB of heap available for socket concurrency.

\begin{table}[htbp]
\caption{Simulation Environment vs. Physical Deployment Realities}
\begin{center}
\begin{tabularx}{\columnwidth}{p{2.0cm}XX}
\toprule
\textbf{Parameter} & \textbf{Wokwi Simulation} & \textbf{Physical Hardware} \\
\midrule
Storage Medium & Virtual FAT32 disk image capped at 256 MB. & Physical MicroSD (4 GB to 32 GB FAT32). \\
\addlinespace
Network Path & Bridged via \texttt{wokwigw} (\texttt{0.0.0.0:9080}). & Direct 2.4 GHz 802.11 b/g/n Wi-Fi or SoftAP. \\
\addlinespace
Current Draw & Virtualized (0 mA). & 80--250 mA typical; 500 mA peak (requires decoupling). \\
\addlinespace
SPI Clocking & Fixed @ 4 MHz for virtual bus synchronization. & 4 MHz to 10 MHz breadboard; 20 MHz on PCB traces. \\
\bottomrule
\end{tabularx}
\label{tab:comparison}
\end{center}
\end{table}

\section{\textbf{Socio-Economic Impact in Bangladesh}}
The deployment of sub-1-Watt, \$7 USD air-gapped storage nodes presents critical practical applications across Bangladesh:
\begin{itemize}
    \item \textbf{Monsoon Flood Relief:} During catastrophic seasonal flooding, cellular basestations frequently fail due to power grid collapse. Battery-powered ESP32 nodes deployed in rural shelters can continuously broadcast survival guides, medical dosages, and evacuation maps over local Wi-Fi.
    \item \textbf{Agricultural Telemetry Logging:} Remote solar-powered agricultural stations monitoring soil moisture and saline intrusion in coastal districts can log sensor data to local SD cards, allowing field engineers to retrieve months of logs wirelessly without physical cable hookups.
    \item \textbf{Off-Grid Rural Education:} Digital textbooks and encyclopedias can be hosted in village schools lacking internet access for a total capital expenditure under 750 BDT, completely bypassing recurring cloud hosting fees.
\end{itemize}

\section{\textbf{Conclusion \& Future Roadmap}}
This project successfully designed, implemented, and empirically verified an autonomous, low-cost embedded NAS and static web host powered by the ESP32 microcontroller. By combining FreeRTOS dual-core task isolation, a 2~KB streaming multipart parser, Flash ROM template embedding, and multi-interface gateway bridging, the node maintains steady throughput across diverse local networking environments without risking memory exhaustion.

Future enhancements include:
\begin{enumerate}
    \item \textbf{Standalone SoftAP Mode (\texttt{192.168.4.1}):} Broadcasting an independent SSID with an internal DHCP daemon for 100\% routerless off-grid operation.
    \item \textbf{ESP32-S3 USB Host Integration:} Migrating to the ESP32-S3 microcontroller to interface commercial USB flash drives via USB Mass Storage Class (MSC).
    \item \textbf{Integrated Battery Management:} Pairing the module with a 18650 lithium battery enclosure for over 20 hours of continuous field runtime.
\end{enumerate}

\begin{thebibliography}{00}
\bibitem{b1} Espressif Systems, \textit{ESP32 Technical Reference Manual}, Version 4.9, Espressif Systems, Shanghai, China, 2023.
\bibitem{b2} R. Barry, \textit{Mastering the FreeRTOS Real Time Kernel: A Hands-On Tutorial Guide}, Real Time Engineers Ltd., 2016.
\bibitem{b3} SD Association, \textit{SD Specifications Part 1: Physical Layer Simplified Specification}, Version 8.00, SD Card Association, San Ramon, CA, USA, 2020.
\bibitem{b4} Wokwi, \textit{Wokwi IoT Simulation Architecture and Virtual Network Bridging}, Documentation Reference, 2024. [Online]. Available: \url{https://docs.wokwi.com}
\bibitem{b5} G. S. Murady, J. Ferdus, and S. Islam, ``Project.MM: ESP32 Embedded NAS \& Static Web Server Workspace,'' GitHub Repository, 2026. [Online]. Available: \url{https://github.com/b1tranger/Microprocessor_Assembly_Lab}
\end{thebibliography}

\appendices

\section{Tri-Platform Architectural Comparison Matrix}
Table~\ref{tab:triplatform} provides a comparative analysis contrasting Wokwi Local simulation, Wokwi Cloud, and physical ESP32 deployment.

\begin{table*}[t]
\caption{Tri-Platform Architectural Comparison Matrix}
\begin{center}
\begin{tabularx}{\textwidth}{p{3.2cm}XXX}
\toprule
\textbf{Architectural Metric} & \textbf{Wokwi Local (VS Code)} & \textbf{Wokwi Cloud (Browser)} & \textbf{Physical Hardware} \\
\midrule
\textbf{Network Accessibility} & Real LAN \& Mobile Hotspot via \texttt{wokwigw.exe} (\texttt{0.0.0.0:9080}). & Restricted virtual private gateway (Requires paid subscription). & Native 802.11 b/g/n Wi-Fi \& Autonomous SoftAP (\texttt{192.168.4.1}). \\
\addlinespace
\textbf{File Upload Execution} & Fully functional streaming uploads committed to host SSD. & Blocked by cloud server upload quotas. & Physical NAND flash block writes via VSPI. \\
\addlinespace
\textbf{Storage Capacity Limit} & Capped at 256 MB by virtual FAT32 image. & Capped at 32 MB in virtual memory. & 4 GB to 32 GB FAT32 MicroSD cards. \\
\addlinespace
\textbf{Power Consumption} & 0 Watts (Host CPU cycles only). & 0 Watts (Remote server computation). & $< 1$ Watt ($< 500$ mW typical operation). \\
\addlinespace
\textbf{Toolchain \& Build} & PlatformIO Core CLI (offline, deterministic). & WebAssembly cloud compiler (internet-dependent). & Flashed via PlatformIO or esptool.py over UART. \\
\bottomrule
\end{tabularx}
\label{tab:triplatform}
\end{center}
\end{table*}

\section{Evaluator Q\&A Defense Quick-Reference}
Below are key technical rebuttals formulated for the project oral defense:
\begin{itemize}
    \item \textbf{Q: Why choose ESP32 over Raspberry Pi?}  
    \textit{Rebuttal:} Raspberry Pi costs \$45--\$75 and consumes 3--7 Watts under load with complex Linux security risks. The ESP32 costs under \$4, consumes $< 0.5$ Watts, boots in under 500~ms, and runs bare-metal FreeRTOS with zero OS vulnerabilities.
    \item \textbf{Q: How does the system handle uploads larger than available RAM?}  
    \textit{Rebuttal:} The firmware never buffers the whole file in RAM. It implements a streaming multipart parser that receives 1.4~KB TCP chunks and immediately writes 512-byte blocks across the VSPI bus to the SD card before acknowledging subsequent packets. RAM usage remains locked at ~2~KB.
    \item \textbf{Q: Is the HTTP server running on the PC or ESP32?}  
    \textit{Rebuttal:} The server is 100\% internal and native to the ESP32 C++ firmware (\texttt{WebServer server(80)}). The PC runs no web server software; \texttt{wokwigw.exe} is merely a transparent Layer 2/3 packet bridge forwarding host port 9080 to the virtual ESP32.
    \item \textbf{Q: What if campus Wi-Fi has AP Client Isolation enabled?}  
    \textit{Rebuttal:} The host laptop activates Windows Mobile Hotspot, creating an independent wireless subnet on deterministic IP \texttt{192.168.137.1}. Client devices connect directly to this hotspot, completely bypassing campus network isolation.
\end{itemize}

\section{Repository Architecture \& Component Responsibilities}
To maintain standalone modularity and allow independent execution across local simulations, cloud browsers, and physical hardware, the workspace decouples toolchains, assets, and firmware. The complete structural tree and file responsibility matrix is summarized below:

\begin{lstlisting}
Project.MM/
|-- PRESENTATION_GUIDE.md    # Presentation & defense guide
|-- README.md                # Project overview & quick links
|-- Presentation/            # Slides, bridges & speaker aids
|   |-- PRESENTATION_GUIDE.md
|   |-- v1_presentationBridge.md
|   `-- v2_presentationBridge.md
|-- FinalReport/             # IEEE report sources & PDFs
|   |-- v4_MM_Lab_Final_Report(PresentationAid).tex
|   `-- v4_MM_Lab_Final_Report(PresentationAid).pdf
|-- LocalConfig/             # Interactive Local Workspace
|   |-- platformio.ini       # PlatformIO ESP32 config
|   |-- wokwi.toml           # 256 MB virtual SD mount
|   |-- diagram.json         # Circuit schematic & SPI wiring
|   |-- wokwigw.exe          # Layer-2/3 TCP gateway binary
|   |-- setup_firewall.bat   # Windows Firewall rule installer
|   |-- refresh_firmware.bat # Firmware refresh launcher
|   |-- refresh_firmware.ps1 # Asset sync & build script
|   |-- build_firmware.bat   # Root build launcher alias
|   |-- runGuide.md          # Setup & 4 network scenarios
|   |-- README.md            # Workspace operational guide
|   |-- .gitignore           # Ignores .pio build outputs
|   |-- data/                # Initial assets & virtual SD
|   |   |-- default/         # Protected files (lab_report.txt)
|   |   |-- www/             # Web UI assets (HTML/CSS)
|   |   `-- network/         # Mapped simulated SD card
|   |       |-- default/     # Seeded default files
|   |       |-- www/         # Hosted static website files
|   |       `-- uploads/     # Target folder for file uploads
|   |-- scripts/             # Asset conversion pipeline
|   |   |-- generate_assets.bat
|   |   |-- generate_assets.ps1 # data/ -> default_assets.h
|   |   `-- generate_assets.py  # PlatformIO pre-build hook
|   |-- include/
|   |   |-- default_assets.h # PROGMEM arrays in Flash ROM
|   |   `-- html_templates.h # CSS styling & HTML templates
|   |-- setup/               # Toolchain provisioning scripts
|   |   |-- SETUP_GUIDE.md
|   |   |-- install_dependencies.ps1
|   |   |-- build_firmware.ps1
|   |   `-- start_gateway.ps1
|   `-- src/
|       `-- main.cpp         # Complete ESP32 firmware
|-- CloudConfig/             # Standalone Wokwi Web Workspace
|   |-- README.md            # Wokwi Web execution guide
|   |-- sketch.ino           # In-simulation loopback client
|   |-- html_templates.h     # Glassmorphic UI templates
|   |-- diagram.json         # Cloud circuit schematic
|   `-- libraries.txt        # Library manifest
`-- doc/                     # Academic documentation & logs
    |-- CLOUD_VS_VSCODE_VS_HARDWARE.md # Tri-platform matrix
    |-- note/                # Assembly & architecture theory
    |-- prompts/             # Session archives & summaries
    `-- research/            # Feasibility & lab proposals
\end{lstlisting}

\section{Complete Firmware Source Code (\texttt{src/main.cpp})}
The full 827-line C++ firmware source code executing on the ESP32 is listed below:

\begin{lstlisting}[language=C++]
#include "default_assets.h"
#include "html_templates.h"
#include <Arduino.h>
#include <HTTPClient.h>
#include <SD.h>
#include <SPI.h>
#include <WebServer.h>
#include <WiFi.h>
#include <vector>

// Hardware SPI Chip Select pin mapped to GPIO 5
#define SD_CS_PIN 5

// Wokwi virtual Wi-Fi access point credentials
const char *ssid = "Wokwi-GUEST";
const char *password = "";

WebServer server(80);
File uploadFile;

// --- Unit Formatting Helper ---
String formatBytes(size_t bytes) {
  if (bytes < 1024)
    return String(bytes) + " B";
  if (bytes < (1024 * 1024))
    return String(bytes / 1024.0, 1) + " KB";
  return String(bytes / (1024.0 * 1024.0), 1) + " MB";
}

// ============================================================================
// SECURITY & FILE PROTECTION SYSTEM
// ============================================================================
bool isProtectedPath(const String &rawPath) {
  String p = rawPath;
  if (!p.startsWith("/"))
    p = "/" + p;

  // Direct default root mirrors or folders
  if (p == "/lab_report.txt")
    return true;
  if (p.startsWith("/default/") || p == "/default")
    return true;
  if (p.startsWith("/www/") || p == "/www")
    return true;

  // Match against all registered embedded assets
  for (size_t i = 0; i < EMBEDDED_ASSET_COUNT; i++) {
    if (p == EMBEDDED_ASSETS[i].sdPath && EMBEDDED_ASSETS[i].isProtected) {
      return true;
    }
  }
  return false;
}

void ensureParentDir(const String &path) {
  int lastSlash = path.lastIndexOf('/');
  if (lastSlash > 0) {
    String dir = path.substring(0, lastSlash);
    if (!SD.exists(dir)) {
      SD.mkdir(dir);
    }
  }
}

// ============================================================================
// PHASE 1: STORAGE INITIALIZATION & DYNAMIC ASSET SEEDING
// ============================================================================
void phase1_seedInitialFiles() {
  // Ensure primary folder structure exists
  if (!SD.exists("/default"))
    SD.mkdir("/default");
  if (!SD.exists("/www"))
    SD.mkdir("/www");
  if (!SD.exists("/uploads"))
    SD.mkdir("/uploads");

  // Dynamically seed all embedded assets from data/ folder
  for (size_t i = 0; i < EMBEDDED_ASSET_COUNT; i++) {
    const EmbeddedAsset &asset = EMBEDDED_ASSETS[i];
    ensureParentDir(asset.sdPath);

    bool needsSeed = !SD.exists(asset.sdPath);
    if (!needsSeed) {
      File checkF = SD.open(asset.sdPath, FILE_READ);
      if (checkF) {
        if (checkF.size() == 0 || checkF.size() != asset.size)
          needsSeed = true;
        checkF.close();
      }
    }

    if (needsSeed) {
      if (SD.exists(asset.sdPath)) {
        SD.remove(asset.sdPath);
      }
      File f = SD.open(asset.sdPath, FILE_WRITE);
      if (f) {
        // Chunk writes into 512-byte blocks to ensure single-sector CMD24
        // execution
        size_t writtenTotal = 0;
        while (writtenTotal < asset.size) {
          size_t chunk = asset.size - writtenTotal;
          if (chunk > 512)
            chunk = 512;
          f.write(asset.data + writtenTotal, chunk);
          writtenTotal += chunk;
        }
        f.flush();
        f.close();
        Serial.printf("[STORAGE] Seeded %s (%u bytes)\n", asset.sdPath,
                      (unsigned int)asset.size);
      }
    }
  }

  // Backwards-compatible mirror for root /lab_report.txt if missing
  if (!SD.exists("/lab_report.txt") && SD.exists("/default/lab_report.txt")) {
    File src = SD.open("/default/lab_report.txt", FILE_READ);
    File dst = SD.open("/lab_report.txt", FILE_WRITE);
    if (src && dst) {
      uint8_t buf[256];
      while (src.available()) {
        size_t len = src.read(buf, sizeof(buf));
        dst.write(buf, len);
      }
      src.close();
      dst.close();
      Serial.println("[STORAGE] Mirrored /default/lab_report.txt to "
                     "/lab_report.txt for compatibility");
    }
  }
}

bool phase1_initStorage() {
  Serial.println("\n--------------------------------------------------");
  Serial.println("  [PHASE 1 EXECUTION] Mounting MicroSD via VSPI   ");
  Serial.println("--------------------------------------------------");

  // Explicitly initialize hardware VSPI bus
  SPI.begin(18, 19, 23, SD_CS_PIN);

  // Initialize SD with reliable 4 MHz SPI clock (standard for VSPI & Wokwi
  // simulation)
  if (!SD.begin(SD_CS_PIN, SPI, 4000000)) {
    Serial.println(
        "\r\n[PHASE 1 ERROR] SD Card initialization failed! Check SPI pins.");
    return false;
  }

  uint64_t totalBytes = SD.totalBytes();
  uint64_t usedBytes = SD.usedBytes();
  Serial.printf("[PHASE 1 SUCCESS] FAT32 Storage Mounted.\n");
  Serial.printf("[PHASE 1 CAPACITY] Total: %llu MB | Used: %llu MB\n",
                totalBytes / (1024 * 1024), usedBytes / (1024 * 1024));

  phase1_seedInitialFiles();

  Serial.println("[PHASE 1 DIRECTORY TREE]:");
  File root = SD.open("/");
  File file = root.openNextFile();
  while (file) {
    if (file.isDirectory()) {
      Serial.printf("  [DIR]  /%s\n", file.name());
    } else {
      Serial.printf("  [FILE] /%s (%s)\n", file.name(),
                    formatBytes(file.size()).c_str());
    }
    file = root.openNextFile();
  }
  root.close();

  Serial.println("--------------------------------------------------\n");
  return true;
}

// ============================================================================
// PHASE 2: HTTP ROUTE HANDLERS
// ============================================================================

struct StorageItem {
  String path;
  size_t size;
  bool isProtected;
};

void collectFilesRecursively(File dir, std::vector<StorageItem> &items,
                             int depth = 0) {
  if (!dir || !dir.isDirectory() || depth > 3)
    return;

  File file = dir.openNextFile();
  while (file) {
    if (file.isDirectory()) {
      collectFilesRecursively(file, items, depth + 1);
    } else {
      String fullPath = String(file.name());
      if (!fullPath.startsWith("/"))
        fullPath = "/" + fullPath;

      // Skip root /lab_report.txt if /default/lab_report.txt exists to avoid
      // duplicate rows
      if (fullPath == "/lab_report.txt" &&
          SD.exists("/default/lab_report.txt")) {
        file = dir.openNextFile();
        continue;
      }

      StorageItem item;
      item.path = fullPath;
      item.size = file.size();
      item.isProtected = isProtectedPath(fullPath);
      items.push_back(item);
    }
    file = dir.openNextFile();
  }
}

// High-performance streaming of SD card files to client in 1024-byte block chunks.
// Eliminates WebServer::streamFile() CPU byte-by-byte bottlenecks and socket deadlocks.
static void streamSDFile(File &file, const String &contentType,
                         bool isAttachment = false,
                         const String &downloadFilename = "") {
  if (!file || file.isDirectory()) {
    if (file)
      file.close();
    server.sendHeader("Connection", "close");
    server.send(404, "text/plain", "File Not Found");
    WiFiClient client = server.client();
    client.flush();
    client.stop();
    return;
  }

  size_t totalSize = file.size();
  file.setBufferSize(1024);

  server.sendHeader("Cache-Control", "no-cache, no-store, must-revalidate");
  server.sendHeader("Pragma", "no-cache");
  server.sendHeader("Expires", "-1");
  if (isAttachment && downloadFilename.length() > 0) {
    server.sendHeader("Content-Disposition",
                      "attachment; filename=\"" + downloadFilename + "\"");
  }
  server.sendHeader("Content-Length", String(totalSize));
  server.sendHeader("Connection", "close");
  server.setContentLength(totalSize);
  server.send(200, contentType, "");

  WiFiClient client = server.client();
  uint8_t buffer[1024];
  size_t remaining = totalSize;

  while (client.connected() && remaining > 0) {
    size_t toRead = (remaining > sizeof(buffer)) ? sizeof(buffer) : remaining;
    size_t nRead = file.read(buffer, toRead);
    if (nRead == 0) {
      break;
    }
    size_t nWritten = client.write(buffer, nRead);
    if (nWritten == 0) {
      break;
    }
    remaining -= nWritten;
    vTaskDelay(1);
  }

  file.close();
  client.flush();
  client.stop();
}

// Route: NAS Dashboard (Streams /www/fileserver.html directly from MicroSD
// storage)
void handleRoot() {
  if (SD.exists("/www/fileserver.html")) {
    File htmlFile = SD.open("/www/fileserver.html", FILE_READ);
    if (htmlFile && htmlFile.size() > 0) {
      streamSDFile(htmlFile, "text/html");
      return;
    }
    if (htmlFile)
      htmlFile.close();
  }

  // Self-healing fallback: re-seed from Flash ROM if missing or 0 bytes
  for (size_t i = 0; i < EMBEDDED_ASSET_COUNT; i++) {
    if (String(EMBEDDED_ASSETS[i].sdPath) == "/www/fileserver.html") {
      ensureParentDir("/www/fileserver.html");
      if (SD.exists("/www/fileserver.html")) {
        SD.remove("/www/fileserver.html");
      }
      File f = SD.open("/www/fileserver.html", FILE_WRITE);
      if (f) {
        size_t writtenTotal = 0;
        while (writtenTotal < EMBEDDED_ASSETS[i].size) {
          size_t chunk = EMBEDDED_ASSETS[i].size - writtenTotal;
          if (chunk > 512)
            chunk = 512;
          f.write(EMBEDDED_ASSETS[i].data + writtenTotal, chunk);
          writtenTotal += chunk;
        }
        f.flush();
        f.close();
        File htmlFile = SD.open("/www/fileserver.html", FILE_READ);
        if (htmlFile && htmlFile.size() > 0) {
          streamSDFile(htmlFile, "text/html");
          return;
        }
        if (htmlFile)
          htmlFile.close();
      }
      // Direct Flash ROM fallback if SD write failed
      server.sendHeader("Connection", "close");
      server.send_P(200, "text/html", (PGM_P)EMBEDDED_ASSETS[i].data,
                    EMBEDDED_ASSETS[i].size);
      WiFiClient client = server.client();
      client.flush();
      client.stop();
      return;
    }
  }

  server.sendHeader("Connection", "close");
  server.send(
      404, "text/plain",
      "File Server Interface Not Found on MicroSD (/www/fileserver.html)");
  WiFiClient client = server.client();
  client.flush();
  client.stop();
}

// Route: API Endpoint to return live storage stats and file directory in JSON
// format
void handleApiFiles() {
  uint64_t total = SD.totalBytes();
  uint64_t used = SD.usedBytes();
  float usedPercent =
      (total > 0) ? ((float)used / (float)total) * 100.0f : 0.0f;

  std::vector<StorageItem> items;
  File root = SD.open("/");
  collectFilesRecursively(root, items);
  root.close();

  String json = "{";
  json += "\"total\":" + String((unsigned long long)total) + ",";
  json += "\"used\":" + String((unsigned long long)used) + ",";
  json += "\"totalFormatted\":\"" + formatBytes(total) + "\",";
  json += "\"usedFormatted\":\"" + formatBytes(used) + "\",";
  json += "\"usedPercent\":" + String(usedPercent, 1) + ",";
  json += "\"files\":[";

  for (size_t i = 0; i < items.size(); i++) {
    const auto &item = items[i];
    int lastSlash = item.path.lastIndexOf('/');
    String dirName = (lastSlash > 0) ? item.path.substring(0, lastSlash) : "/";
    String baseName =
        (lastSlash >= 0) ? item.path.substring(lastSlash + 1) : item.path;

    if (i > 0)
      json += ",";
    json += "{";
    json += "\"path\":\"" + item.path + "\",";
    json += "\"name\":\"" + baseName + "\",";
    json += "\"dir\":\"" + dirName + "\",";
    json += "\"size\":" + String(item.size) + ",";
    json += "\"sizeFormatted\":\"" + formatBytes(item.size) + "\",";
    json += "\"isProtected\":" + String(item.isProtected ? "true" : "false");
    json += "}";
  }
  json += "]}";

  server.sendHeader("Connection", "close");
  server.send(200, "application/json", json);
  WiFiClient client = server.client();
  client.flush();
  client.stop();
}

// Static Upload State Tracker
static String currentUploadTarget = "";
static size_t uploadFileWrittenBytes = 0;

// Route: File Upload Handler (Direct Streaming Multipart Parser restored from
// cf175466)
void handleFileUpload() {
  HTTPUpload &upload = server.upload();

  if (upload.status == UPLOAD_FILE_START) {
    uploadFileWrittenBytes = 0;
    String filename = upload.filename;
    int slashIdx = filename.lastIndexOf('/');
    if (slashIdx >= 0)
      filename = filename.substring(slashIdx + 1);
    int bslashIdx = filename.lastIndexOf('\\');
    if (bslashIdx >= 0)
      filename = filename.substring(bslashIdx + 1);

    // Sanitize punctuation and illegal FAT32 characters
    String cleanFilename = "";
    for (size_t i = 0; i < filename.length(); i++) {
      char c = filename[i];
      if (c == '\'' || c == '\"' || c == '?' || c == '*' || c == ':' ||
          c == '|' || c == '<' || c == '>') {
        cleanFilename += '_';
      } else {
        cleanFilename += c;
      }
    }
    filename = cleanFilename;

    if (!SD.exists("/uploads")) {
      SD.mkdir("/uploads");
    }
    currentUploadTarget = "/uploads/" + filename;
    Serial.printf("[UPLOAD START] Destination: %s\r\n",
                  currentUploadTarget.c_str());

    // Purge any stale existing file
    if (SD.exists(currentUploadTarget)) {
      SD.remove(currentUploadTarget);
    }

    uploadFile = SD.open(currentUploadTarget, FILE_WRITE);
    if (!uploadFile) {
      Serial.printf("[UPLOAD ERROR] Failed to create: %s\r\n",
                    currentUploadTarget.c_str());
    } else {
      uploadFile.setBufferSize(512);
    }
  } else if (upload.status == UPLOAD_FILE_WRITE) {
    if (uploadFile && upload.currentSize > 0) {
      size_t written = uploadFile.write(upload.buf, upload.currentSize);
      uploadFileWrittenBytes += written;
      if (written != upload.currentSize) {
        Serial.printf("[UPLOAD WRITE MISMATCH] Buffer: %u, Wrote: %u\r\n",
                      (unsigned int)upload.currentSize, (unsigned int)written);
      }
    }
  } else if (upload.status == UPLOAD_FILE_END) {
    if (uploadFile) {
      uploadFile.flush();
      size_t finalFileSize = uploadFile.size();
      uploadFile.close();
      Serial.printf("[UPLOAD SUCCESS] Path: %s | Network: %u bytes | SD Committed: %u bytes | File Size: %u bytes\r\n",
                    currentUploadTarget.c_str(),
                    (unsigned int)upload.totalSize,
                    (unsigned int)uploadFileWrittenBytes,
                    (unsigned int)finalFileSize);
    }
    currentUploadTarget = "";
  } else if (upload.status == UPLOAD_FILE_ABORTED) {
    if (uploadFile) {
      uploadFile.close();
      Serial.println("[UPLOAD ABORTED]\r\n");
    }
    if (currentUploadTarget.length() > 0 && SD.exists(currentUploadTarget)) {
      SD.remove(currentUploadTarget);
      Serial.printf("[UPLOAD ABORTED] Purged incomplete file: %s\r\n",
                    currentUploadTarget.c_str());
    }
    currentUploadTarget = "";
  }
}

// Route: File Deletion Handler (Rejects protected default files)
void handleDelete() {
  if (!server.hasArg("file")) {
    server.sendHeader("Connection", "close");
    server.send(400, "text/plain", "Missing file parameter");
    WiFiClient client = server.client();
    client.flush();
    client.stop();
    return;
  }

  String target = server.arg("file");
  if (!target.startsWith("/"))
    target = "/" + target;

  // Block attempts to delete protected system files
  if (isProtectedPath(target)) {
    Serial.printf(
        "[SECURITY BLOCKED] Deletion rejected for protected file: %s\n",
        target.c_str());
    server.sendHeader("Connection", "close");
    server.send(403, "text/plain",
                "Forbidden: Default system files are write-protected and "
                "cannot be deleted.");
    WiFiClient client = server.client();
    client.flush();
    client.stop();
    return;
  }

  bool removed = false;
  if (SD.exists(target)) {
    SD.remove(target);
    removed = true;
    Serial.printf("[STORAGE] Deleted: %s\n", target.c_str());
  } else {
    String uploadPath = "/uploads" + target;
    if (SD.exists(uploadPath)) {
      SD.remove(uploadPath);
      removed = true;
      Serial.printf("[STORAGE] Deleted: %s\n", uploadPath.c_str());
    }
  }

  if (!removed) {
    Serial.printf("[STORAGE WARN] File not found for deletion: %s\n",
                  target.c_str());
  }

  // Redirect back to dashboard
  server.sendHeader("Location", "/");
  server.sendHeader("Connection", "close");
  server.send(303);
  WiFiClient client = server.client();
  client.flush();
  client.stop();
}

// Route: Static Web Host (Streams /www/index.html from SD card)
void handleSite() {
  if (SD.exists("/www/index.html")) {
    File htmlFile = SD.open("/www/index.html", FILE_READ);
    if (htmlFile && htmlFile.size() > 0) {
      streamSDFile(htmlFile, "text/html");
      return;
    }
    if (htmlFile)
      htmlFile.close();
  }

  // Self-healing fallback: re-seed from Flash ROM if missing or 0 bytes
  for (size_t i = 0; i < EMBEDDED_ASSET_COUNT; i++) {
    if (String(EMBEDDED_ASSETS[i].sdPath) == "/www/index.html") {
      ensureParentDir("/www/index.html");
      if (SD.exists("/www/index.html")) {
        SD.remove("/www/index.html");
      }
      File f = SD.open("/www/index.html", FILE_WRITE);
      if (f) {
        size_t writtenTotal = 0;
        while (writtenTotal < EMBEDDED_ASSETS[i].size) {
          size_t chunk = EMBEDDED_ASSETS[i].size - writtenTotal;
          if (chunk > 512)
            chunk = 512;
          f.write(EMBEDDED_ASSETS[i].data + writtenTotal, chunk);
          writtenTotal += chunk;
        }
        f.flush();
        f.close();
        File htmlFile = SD.open("/www/index.html", FILE_READ);
        if (htmlFile && htmlFile.size() > 0) {
          streamSDFile(htmlFile, "text/html");
          return;
        }
        if (htmlFile)
          htmlFile.close();
      }
      // Direct Flash ROM fallback if SD write failed
      server.sendHeader("Connection", "close");
      server.send_P(200, "text/html", (PGM_P)EMBEDDED_ASSETS[i].data,
                    EMBEDDED_ASSETS[i].size);
      WiFiClient client = server.client();
      client.flush();
      client.stop();
      return;
    }
  }

  server.sendHeader("Connection", "close");
  server.send(404, "text/plain",
              "Static Site Not Found on MicroSD (/www/index.html)");
  WiFiClient client = server.client();
  client.flush();
  client.stop();
}

// Route: File Download Stream (High-throughput chunked streaming)
void handleDownload() {
  if (!server.hasArg("file")) {
    server.sendHeader("Connection", "close");
    server.send(400, "text/plain", "Missing file parameter");
    WiFiClient client = server.client();
    client.flush();
    client.stop();
    return;
  }

  String req = server.arg("file");
  if (!req.startsWith("/"))
    req = "/" + req;

  String resolvedPath = "";
  if (SD.exists(req)) {
    resolvedPath = req;
  } else if (SD.exists("/default" + req)) {
    resolvedPath = "/default" + req;
  } else if (SD.exists("/www" + req)) {
    resolvedPath = "/www" + req;
  } else if (SD.exists("/uploads" + req)) {
    resolvedPath = "/uploads" + req;
  }

  if (resolvedPath.length() == 0) {
    server.sendHeader("Connection", "close");
    server.send(404, "text/plain", "File Not Found on Storage");
    WiFiClient client = server.client();
    client.flush();
    client.stop();
    return;
  }

  File downloadFile = SD.open(resolvedPath, FILE_READ);
  if (!downloadFile || downloadFile.isDirectory()) {
    if (downloadFile)
      downloadFile.close();
    server.sendHeader("Connection", "close");
    server.send(404, "text/plain", "File Not Found on Storage");
    WiFiClient client = server.client();
    client.flush();
    client.stop();
    return;
  }

  int lastSlash = resolvedPath.lastIndexOf('/');
  String dlName =
      (lastSlash >= 0) ? resolvedPath.substring(lastSlash + 1) : resolvedPath;

  Serial.printf("[DOWNLOAD] Streaming: %s (%u bytes)\r\n",
                resolvedPath.c_str(), (unsigned int)downloadFile.size());

  streamSDFile(downloadFile, "application/octet-stream", true, dlName);
}

// Helper function for reliable diagnostic requests in virtual simulation
static int performTestRequest(const String &url) {
  HTTPClient http;
  http.setReuse(false);
  http.setTimeout(8000);
  if (!http.begin(url)) {
    return -99;
  }
  int code = http.GET();
  // Fully drain the response stream to ensure graceful socket teardown
  WiFiClient *stream = http.getStreamPtr();
  if (stream) {
    uint8_t dummy[256];
    while (http.connected() && stream->available() > 0) {
      stream->read(dummy, sizeof(dummy));
    }
  }
  http.end();
  // Yield to loopTask on Core 1 to allow WebServer to complete TCP close &
  // rearm
  vTaskDelay(800 / portTICK_PERIOD_MS);
  return code;
}

// Background Self-Diagnostics on Core 1
void runClientTests(void *parameter) {
  vTaskDelay(4000 / portTICK_PERIOD_MS);

  Serial.println("\r\n==================================================");
  Serial.println("  [PHASE 2 AUTO-TEST] Running Network Diagnostics ");
  Serial.println("==================================================");

  // Test 1: Dashboard explorer
  int code = performTestRequest("http://127.0.0.1/");
  Serial.printf("[TEST 1: Root Explorer] Response: %d\r\n", code);

  // Test 2: Static site stream
  code = performTestRequest("http://127.0.0.1/site");
  Serial.printf("[TEST 2: Static Host /site] Response: %d\r\n", code);

  // Test 3: Download default asset
  code = performTestRequest("http://127.0.0.1/get?file=lab_report.txt");
  Serial.printf("[TEST 3: Download /get] Response: %d\r\n", code);

  // Test 4: Verify default files are protected against deletion
  code = performTestRequest(
      "http://127.0.0.1/delete?file=/default/lab_report.txt");
  Serial.printf(
      "[TEST 4: Delete Protection] Response: %d (Expected: 403 Forbidden)\r\n",
      code);

  // Test 5: Verify JSON files API
  code = performTestRequest("http://127.0.0.1/api/files");
  Serial.printf("[TEST 5: API /api/files] Response: %d\r\n", code);

  // Test 6: Verify SD Card file writing (Single 512-byte Sector)
  if (!SD.exists("/uploads")) {
    SD.mkdir("/uploads");
  }
  File testF = SD.open("/uploads/diag_test512.tmp", FILE_WRITE);
  if (testF) {
    testF.setBufferSize(512);
    alignas(4) uint8_t dummy[512];
    memset(dummy, 'A', sizeof(dummy));
    errno = 0;
    size_t w = testF.write(dummy, 512);
    testF.flush();
    testF.close();
    Serial.printf("[TEST 6: SD 512B Block Write] Wrote: %u / 512 bytes (errno: "
                  "%d, %s)\r\n",
                  (unsigned int)w, errno, strerror(errno));
    SD.remove("/uploads/diag_test512.tmp");
  } else {
    Serial.printf(
        "[TEST 6: SD 512B Block Write] Open failed (errno: %d, %s)\r\n", errno,
        strerror(errno));
  }

  // Test 7: Verify SD Card multi-block writeback (2x 512-byte blocks)
  File testF2 = SD.open("/uploads/diag_test1024.tmp", FILE_WRITE);
  if (testF2) {
    testF2.setBufferSize(512);
    alignas(4) uint8_t dummy2[512];
    memset(dummy2, 'B', sizeof(dummy2));
    errno = 0;
    size_t w1 = testF2.write(dummy2, 512);
    size_t w2 = testF2.write(dummy2, 512);
    testF2.flush();
    testF2.close();
    Serial.printf("[TEST 7: SD 1024B Multi-Block Write] Wrote: %u + %u bytes "
                  "(errno: %d, %s)\r\n",
                  (unsigned int)w1, (unsigned int)w2, errno, strerror(errno));
    SD.remove("/uploads/diag_test1024.tmp");
  } else {
    Serial.printf(
        "[TEST 7: SD 1024B Multi-Block Write] Open failed (errno: %d, %s)\r\n",
        errno, strerror(errno));
  }

  Serial.println("==================================================\r\n");
  vTaskDelete(NULL);
}

// ============================================================================
// SYSTEM INITIALIZATION & MAIN LOOP
// ============================================================================
void setup() {
  Serial.begin(115200);
  delay(1000);

  // Mount SD card & seed initial files
  if (!phase1_initStorage()) {
    Serial.println("[ABORT] Cannot proceed without storage.");
    return;
  }

  // Connect to virtual Wi-Fi (explicit STA mode + channel 6 skips slow
  // multi-channel scan in Wokwi)
  WiFi.mode(WIFI_STA);
  Serial.printf("[NET] Connecting to Wi-Fi SSID: %s (ch 6)...", ssid);
  WiFi.begin(ssid, password, 6);
  int attempts = 0;
  while (WiFi.status() != WL_CONNECTED) {
    delay(200);
    Serial.print(".");
    attempts++;
    if (attempts % 50 == 0) {
      Serial.printf("\n[NET WARN] Still waiting for Wi-Fi (status=%d). Ensure "
                    "wokwigw.exe is running on ws://localhost:9011 before "
                    "starting Wokwi.\n",
                    (int)WiFi.status());
      WiFi.disconnect(false);
      WiFi.begin(ssid, password, 6);
    }
  }
  Serial.println("\n[NET] Wi-Fi Connected!");
  Serial.printf("[NET] Node IP Address: %s\n",
                WiFi.localIP().toString().c_str());

  // Bind WebServer Routes
  server.on("/", HTTP_GET, handleRoot);
  server.on("/fileserver.html", HTTP_GET, handleRoot);
  server.on("/www/fileserver.html", HTTP_GET, handleRoot);
  server.on("/api/files", HTTP_GET, handleApiFiles);
  server.on("/site", HTTP_GET, handleSite);
  server.on("/index.html", HTTP_GET, handleSite);
  server.on("/www/index.html", HTTP_GET, handleSite);
  server.on("/get", HTTP_GET, handleDownload);
  server.on("/delete", HTTP_GET, handleDelete);

  // Bind Upload handler: completion callback + streaming upload callback
  server.on(
      "/upload", HTTP_POST,
      []() {
        HTTPUpload &upload = server.upload();
        server.sendHeader("Connection", "close");
        if (upload.status == UPLOAD_FILE_END) {
          server.send(200, "application/json",
                      "{\"status\":\"ok\",\"bytes\":" +
                          String((unsigned long)upload.totalSize) + "}");
        } else {
          server.send(500, "application/json",
                      "{\"status\":\"error\",\"message\":\"Upload aborted or "
                      "incomplete payload received\"}");
        }
        WiFiClient client = server.client();
        client.flush();
        client.stop();
      },
      handleFileUpload);

  // Catch unhandled endpoints (e.g. browser favicon.ico probes) to prevent socket starvation
  server.onNotFound([]() {
    server.sendHeader("Connection", "close");
    server.send(404, "text/plain", "Not Found");
    WiFiClient client = server.client();
    client.flush();
    client.stop();
  });

  server.begin();
  Serial.println(
      "[HTTP] Server listening on port 80 with File Upload enabled.");

  // Launch background diagnostics on Core 1
  xTaskCreatePinnedToCore(runClientTests, "ClientTest", 6144, NULL, 1, NULL, 1);
}

void loop() { server.handleClient(); }

\end{lstlisting}

\section{Minimalist Industrial NAS Interface (\texttt{data/www/fileserver.html})}
The responsive, dark-mode file manager interface streamed from the MicroSD storage is listed below (CSS styling in the header is truncated for brevity):

\begin{lstlisting}[language=html]
<!DOCTYPE html>
<html lang="en">

<head>
  <meta charset="UTF-8">
  <meta name="viewport" content="width=device-width, initial-scale=1.0">
  <title>ESP32 Cloud NAS & Storage Node</title>
  <link rel="stylesheet"
    href="https://fonts.googleapis.com/css2?family=Inter:wght@300;400;500;600;700">
  <link rel="preconnect" href="https://fonts.gstatic.com" crossorigin>
  <link rel="stylesheet"
    href="https://fonts.googleapis.com/css2?family=Inter:wght@300;400;500;600;700">
  <style>
    /* ... [CSS Minimalist Industrial Dark-Theme Styles Truncated for Brevity] ... */
  </style>
</head>

<body>
  <div class="container">
    <div class="header-card">
      <div class="title-group">
        <h1>ESP32 Storage NAS Node</h1>
        <p>Hardware VSPI FAT32 MicroSD Storage Engine</p>
      </div>
      <div class="header-actions">
        <span id="syncIndicator" class="sync-badge" title="Auto-refreshing file list every 10 seconds (Click to refresh now)">
          <span class="sync-dot"></span> [sync] Auto-sync: <span id="countdown">10</span>s
        </span>
        <a href="index.html" class="btn-site">[web] Static Site</a>
      </div>
    </div>

    <!-- Storage Meter -->
    <div class="meter-card" id="meterSection">
      <div class="meter-header">
        <span class="meter-title">Storage Utilization</span>
        <span class="meter-stats" id="storageSummary">Loading...</span>
      </div>
      <div class="progress-bar-bg">
        <div class="progress-bar-fill" id="meterBarFill" style="width: 0%;"></div>
      </div>
    </div>

    <!-- Upload Box -->
    <div class="upload-box">
      <form id="uploadForm" class="upload-form" onsubmit="event.preventDefault(); submitUpload();">
        <input type="file" id="fileSelect" name="upload_file" required>
        <div id="fileDetailCard" class="file-detail-card" style="display: none;">
          <div class="file-detail-main">
            <span class="file-detail-icon">[doc]</span>
            <div class="file-detail-meta">
              <div class="file-detail-name" id="chosenFileName">No file selected</div>
              <div class="file-detail-tags">
                <span class="file-detail-badge" id="chosenFileSize">0 B</span>
                <span class="file-detail-badge" id="chosenFileType">Unknown</span>
                <span class="file-detail-status">Ready to Upload</span>
              </div>
            </div>
          </div>
          <button type="button" class="btn-clear-file" id="clearFileBtn" title="Clear selection" onclick="clearFileSelection();">[x]</button>
        </div>
        <button type="submit" id="uploadBtn" class="btn-upload">[upload] Upload to Storage</button>
      </form>
      <div id="progressBox" class="progress-box">
        <div class="progress-info">
          <span id="fileNameDisplay"><strong>Uploading:</strong> </span>
          <span id="percentDisplay" style="color: var(--accent); font-weight: 700;">0%</span>
        </div>
        <div class="progress-bar-bg">
          <div id="progressBarFill" class="progress-bar-fill" style="width: 0%;"></div>
        </div>
        <div class="progress-info" style="margin-top: 10px; margin-bottom: 0;">
          <span id="bytesDisplay" style="color: var(--text-muted);">0 B / 0 B</span>
          <span id="speedDisplay" class="speed-badge">[speed] 0 KB/s</span>
          <span id="etaDisplay" class="eta-badge">[time] Calculating...</span>
        </div>
      </div>
    </div>

    <!-- File Table -->
    <div class="table-card">
      <table id="fileTable">
        <thead>
          <tr>
            <th>File Name & Location</th>
            <th>Size</th>
            <th style="text-align: right;">Actions</th>
          </tr>
        </thead>
        <tbody id="fileTableBody">
          <!-- Populated dynamically or server-side -->
          <tr>
            <td><strong>[doc] lab_report.txt</strong> <span class="tag-folder">/default</span></td>
            <td>778 B</td>
            <td style="text-align: right;">
              <a href="/get?file=/default/lab_report.txt" class="btn-action btn-dl" title="Download lab_report.txt" aria-label="Download"><span class="btn-icon">[download]</span><span class="btn-text"> Download</span></a>
              <span class="badge-protected" title="Protected System File" aria-label="Protected"><span class="btn-icon">[lock]</span><span class="btn-text"> Protected</span></span>
            </td>
          </tr>
          <tr>
            <td><strong>[doc] index.html</strong> <span class="tag-folder">/www</span></td>
            <td>3.2 KB</td>
            <td style="text-align: right;">
              <a href="/get?file=/www/index.html" class="btn-action btn-dl" title="Download index.html" aria-label="Download"><span class="btn-icon">[download]</span><span class="btn-text"> Download</span></a>
              <span class="badge-protected" title="Protected System File" aria-label="Protected"><span class="btn-icon">[lock]</span><span class="btn-text"> Protected</span></span>
            </td>
          </tr>
        </tbody>
      </table>
    </div>
  </div>

  <script>
    function fmtBytes(b) {
      if (b < 1024) return b + ' B';
      if (b < 1048576) return (b / 1024).toFixed(1) + ' KB';
      return (b / 1048576).toFixed(1) + ' MB';
    }

    function escapeHtml(str) {
      if (!str) return '';
      return String(str)
        .replace(/&/g, '&amp;')
        .replace(/</g, '&lt;')
        .replace(/>/g, '&gt;')
        .replace(/"/g, '&quot;')
        .replace(/'/g, '&#039;');
    }

    function loadFiles() {
      fetch('/api/files')
        .then(function (res) {
          if (!res.ok) throw new Error('API returned status: ' + res.status);
          return res.json();
        })
        .then(function (data) {
          var summary = document.getElementById('storageSummary');
          var fill = document.getElementById('meterBarFill');
          if (summary && data.usedFormatted && data.totalFormatted) {
            summary.textContent = data.usedFormatted + ' / ' + data.totalFormatted + ' (' + data.usedPercent + '% Used)';
          }
          if (fill && typeof data.usedPercent !== 'undefined') {
            fill.style.width = data.usedPercent + '%';
          }

          var tbody = document.getElementById('fileTableBody');
          if (!tbody) return;
          tbody.innerHTML = '';

          if (!data.files || data.files.length === 0) {
            tbody.innerHTML = '<tr><td colspan="3" style="text-align: center; color: var(--text-dim); padding: 24px;">No files found on storage. Upload one above!</td></tr>';
            return;
          }

          data.files.forEach(function (item) {
            var tr = document.createElement('tr');

            var nameTd = document.createElement('td');
            nameTd.innerHTML = '<strong>[doc] ' + escapeHtml(item.name) + '</strong> <span class="tag-folder">' + escapeHtml(item.dir) + '</span>';

            var sizeTd = document.createElement('td');
            sizeTd.textContent = item.sizeFormatted || fmtBytes(item.size);

            var actTd = document.createElement('td');
            actTd.style.textAlign = 'right';

            var dlBtn = document.createElement('a');
            dlBtn.href = '/get?file=' + encodeURIComponent(item.path);
            dlBtn.className = 'btn-action btn-dl';
            dlBtn.title = 'Download ' + item.name;
            dlBtn.setAttribute('aria-label', 'Download ' + item.name);
            dlBtn.innerHTML = '<span class="btn-icon">[download]</span><span class="btn-text"> Download</span>';
            actTd.appendChild(dlBtn);

            if (item.isProtected) {
              var protSpan = document.createElement('span');
              protSpan.className = 'badge-protected';
              protSpan.title = 'Protected System File';
              protSpan.setAttribute('aria-label', 'Protected File');
              protSpan.innerHTML = '<span class="btn-icon">[lock]</span><span class="btn-text"> Protected</span>';
              actTd.appendChild(protSpan);
            } else {
              var delBtn = document.createElement('a');
              delBtn.href = '/delete?file=' + encodeURIComponent(item.path);
              delBtn.className = 'btn-action btn-del';
              delBtn.title = 'Delete ' + item.name;
              delBtn.setAttribute('aria-label', 'Delete ' + item.name);
              delBtn.innerHTML = '<span class="btn-icon">[del]</span><span class="btn-text"> Delete</span>';
              delBtn.onclick = function () {
                return confirm('Delete ' + item.name + '?');
              };
              actTd.appendChild(delBtn);
            }

            tr.appendChild(nameTd);
            tr.appendChild(sizeTd);
            tr.appendChild(actTd);
            tbody.appendChild(tr);
          });
        })
        .catch(function (err) {
          // Keep static preview rows if running in standalone/offline file:/// mode
          console.log('Running in local/offline preview mode:', err.message);
        });
    }

    document.addEventListener('DOMContentLoaded', loadFiles);

    // --- Auto Page Reload & File Sync Timer (10 Seconds) ---
    var countdown = 10;
    var isUploading = false;

    function tickAutoSync() {
      var cdElem = document.getElementById('countdown');
      if (isUploading) {
        if (cdElem) cdElem.textContent = '--';
        return;
      }
      countdown--;
      if (countdown <= 0) {
        countdown = 10;
        loadFiles();
      }
      if (cdElem) cdElem.textContent = countdown;
    }

    var autoSyncTimer = setInterval(tickAutoSync, 1000);

    var syncBadge = document.getElementById('syncIndicator');
    if (syncBadge) {
      syncBadge.onclick = function () {
        if (!isUploading) {
          countdown = 10;
          var cdElem = document.getElementById('countdown');
          if (cdElem) cdElem.textContent = '10';
          loadFiles();
        }
      };
    }

    // Dynamic File Details Display on Selection (Mobile & Desktop)
    var fileSelectInput = document.getElementById('fileSelect');
    if (fileSelectInput) {
      fileSelectInput.addEventListener('change', function () {
        var card = document.getElementById('fileDetailCard');
        if (this.files && this.files.length > 0) {
          var f = this.files[0];
          document.getElementById('chosenFileName').textContent = f.name;
          document.getElementById('chosenFileSize').textContent = fmtBytes(f.size);
          var ext = f.name.indexOf('.') !== -1 ? f.name.split('.').pop().toUpperCase() : 'FILE';
          document.getElementById('chosenFileType').textContent = f.type ? f.type : (ext + ' Document');
          card.style.display = 'flex';
        } else {
          card.style.display = 'none';
        }
      });
    }

    function clearFileSelection() {
      var inp = document.getElementById('fileSelect');
      if (inp) inp.value = '';
      var card = document.getElementById('fileDetailCard');
      if (card) card.style.display = 'none';
    }

    function submitUpload() {
      var inp = document.getElementById('fileSelect');
      if (!inp.files.length) return;
      var f = inp.files[0];

      var box = document.getElementById('progressBox');
      var bar = document.getElementById('progressBarFill');
      var pct = document.getElementById('percentDisplay');
      var name = document.getElementById('fileNameDisplay');
      var bytes = document.getElementById('bytesDisplay');
      var speed = document.getElementById('speedDisplay');
      var eta = document.getElementById('etaDisplay');
      var btn = document.getElementById('uploadBtn');

      isUploading = true;
      box.style.display = 'block';
      name.innerHTML = '<strong>Uploading:</strong> ' + f.name;
      btn.disabled = true;
      btn.style.opacity = '0.5';

      var fd = new FormData();
      fd.append('upload_file', f);

      var xhr = new XMLHttpRequest();
      var t0 = Date.now();
      var lastT = t0;
      var lastL = 0;
      var smoothedSpeed = 0;

      xhr.upload.onprogress = function (e) {
        if (e.lengthComputable) {
          var now = Date.now();
          var p = Math.round((e.loaded / e.total) * 100);
          bar.style.width = p + '%';
          pct.textContent = p + '%';
          bytes.textContent = fmtBytes(e.loaded) + ' / ' + fmtBytes(e.total);

          var dt = (now - lastT) / 1000;
          if (dt >= 0.2 || e.loaded === e.total) {
            var inst = (e.loaded - lastL) / (dt || 0.001);
            smoothedSpeed = (smoothedSpeed === 0) ? inst : (0.65 * smoothedSpeed + 0.35 * inst);
            speed.innerHTML = '[speed] ' + fmtBytes(smoothedSpeed) + '/s';

            var rem = e.total - e.loaded;
            if (p === 100) {
              name.innerHTML = '<strong>Writing to Storage:</strong> ' + f.name;
              speed.innerHTML = '[speed] Committing...';
              eta.textContent = '[time] Finalizing SD card write';
            } else if (smoothedSpeed > 0) {
              var s = Math.round(rem / smoothedSpeed);
              eta.textContent = '[time] ' + (s >= 60 ? Math.floor(s / 60) + 'm ' + (s % 60) + 's' : s + 's left');
            }
            lastT = now;
            lastL = e.loaded;
          }
        }
      };

      xhr.onload = function () {
        isUploading = false;
        countdown = 10;
        if (xhr.status >= 200 && xhr.status < 400) {
          bar.style.width = '100%';
          pct.textContent = '100%';
          name.innerHTML = '<strong style="color: var(--success);">[ok] Upload Complete!</strong> Stored on host SSD/HDD.';
          speed.innerHTML = '[speed] Complete';
          eta.textContent = '[time] Done';
          btn.disabled = false;
          btn.style.opacity = '1';
          inp.value = '';
          clearFileSelection();
          loadFiles();
          setTimeout(function () {
            box.style.display = 'none';
          }, 3000);
        } else {
          var errMsg = 'Upload Error (' + xhr.status + ')';
          try {
            var res = JSON.parse(xhr.responseText);
            if (res.message) errMsg = res.message;
          } catch (e) {}
          name.innerHTML = '<strong style="color: var(--danger);">[err] ' + errMsg + '</strong>';
          btn.disabled = false;
          btn.style.opacity = '1';
        }
      };

      xhr.onerror = function () {
        isUploading = false;
        countdown = 10;
        name.innerHTML = '<strong style="color: var(--danger);">[err] Network Transmission Failed</strong>';
        btn.disabled = false;
        btn.style.opacity = '1';
      };

      xhr.open('POST', '/upload', true);
      xhr.send(fd);
    }
  </script>
</body>

</html>

\end{lstlisting}

\section{Hosted Static Web Document (\texttt{data/www/index.html})}
The static landing page hosted directly on raw MicroSD sectors and delivered via the \texttt{/site} route is listed below (CSS styling in the header is truncated for brevity):

\begin{lstlisting}[language=html]
<!DOCTYPE html>
<html lang="en">

<head>
  <meta charset="UTF-8">
  <meta name="viewport" content="width=device-width, initial-scale=1.0">
  <title>ESP32 Static Edge Host</title>
  <link rel="stylesheet"
    href="https://fonts.googleapis.com/css2?family=Inter:wght@300;400;500;600;700">
  <link rel="preconnect" href="https://fonts.gstatic.com" crossorigin>
  <link rel="stylesheet"
    href="https://fonts.googleapis.com/css2?family=Inter:wght@300;400;500;600;700">
  <style>
    /* ... [CSS Responsive Layout Styles Truncated for Brevity] ... */
  </style>
</head>

<body>
  <div class="card">
    <div class="tag">STATIC EDGE NODE</div>
    <h1>Hosted on MicroSD Storage</h1>
    <p class="lead">
      This static web document is stored as raw sectors on the virtual FAT32 MicroSD card and streamed directly to your
      browser over hardware SPI.
    </p>
    <div class="specs-grid">
      <div class="spec-item">
        <div class="spec-label">Microcontroller</div>
        <div class="spec-val">ESP32 DevKit V1</div>
      </div>
      <div class="spec-item">
        <div class="spec-label">Storage Bus</div>
        <div class="spec-val">Hardware VSPI</div>
      </div>
      <div class="spec-item">
        <div class="spec-label">File Protocol</div>
        <div class="spec-val">FAT32 Block I/O</div>
      </div>
      <div class="spec-item">
        <div class="spec-label">Web Engine</div>
        <div class="spec-val">ESP32 WebServer:80</div>
      </div>
    </div>
    <div class="btn-group">
      <a href="fileserver.html" class="btn">[dir] Go to NAS File Browser</a>
    </div>
  </div>
</body>

</html>
\end{lstlisting}

\end{document}